\section{嘉宾介绍}

翁家翌（Jiayi Weng），2016年入读清华大学计算机系本科，2020年赴卡内基梅隆大学（CMU）攻读硕士，2022年加入 OpenAI。他是 OpenAI 内部 \textbf{Post-Training RL Infrastructure} 的核心搭建者——从 ChatGPT（GPT-3.5）、GPT-4o 到 GPT-5，OpenAI 发布的每一个大模型背后都有他的名字。

在加入 OpenAI 之前，翁家翌已经通过开源项目产生了广泛影响：他在清华开源了全部作业和课程资料以打破信息差；创建了强化学习框架\textbf{天授（Tianshou）}，成为 RL 社区最受欢迎的轻量级框架之一；疫情期间开发了免费签证查询系统 \textbf{tuixue.online}，累计服务数百万次访问。

本期 WhynotTV Podcast 是一次跨越两小时的深度对话，从翁家翌的童年聊起，涵盖求学、开源、职业选择、OpenAI 内部的工作方式，以及他对 AI 未来和人生意义的哲学思考。

\begin{knowledgebox}{访谈背景}
本期播客的 outline 是主持人用 GPT-5 的 Deep Research 功能准备的——而翁家翌本人正是 GPT-5 背后的核心开发者之一。主持人称之为"一个奇妙的闭环"。
\end{knowledgebox}

\subsection{本章小结}

翁家翌的三个关键词是：\textbf{强化学习、Post-Training、Infra}。但他的故事远不止技术本身——他是一个用代码做"慈善"的人，相信开源和信息平权，并在世界 AI 风暴的中心保持着独立的思考。


